% ============================================================
%  CHƯƠNG 6: THỰC NGHIỆM VÀ KẾT QUẢ
%  Người viết: [Họ tên thành viên phụ trách]
% ============================================================

\section{Thực nghiệm và Kết quả}
\label{sec:experiments}

% -----------------------------------------------------------
\subsection{Môi trường thực nghiệm}
% -----------------------------------------------------------
% Gợi ý nội dung:
%   - Ngôn ngữ lập trình, thư viện (Python, scikit-learn, NumPy, …)
%   - Cấu hình phần cứng (CPU, RAM)
%   - Mã nguồn được lưu tại: code/ (xem thư mục code/ trong repo)

% TODO: Viết nội dung ở đây


% -----------------------------------------------------------
\subsection{Dữ liệu tổng hợp}
% -----------------------------------------------------------
% Gợi ý nội dung:
%   - Mô tả cách sinh dữ liệu: make_blobs / make_classification / tự sinh
%   - Số mẫu N, số chiều D, số cụm thực K_true
%   - Các trường hợp: cụm tròn, cụm ellipsoid, cụm chồng lên nhau

\subsubsection{Thiết lập thí nghiệm}
% TODO: Viết nội dung ở đây

\subsubsection{Kết quả phân cụm}
% TODO: Mô tả và phân tích kết quả

% Ví dụ chèn hình kết quả:
% \begin{figure}[H]
%   \centering
%   \begin{subfigure}{0.48\linewidth}
%     \includegraphics[width=\linewidth]{figures/gmm_synthetic_result.pdf}
%     \caption{GMM-EM}
%   \end{subfigure}
%   \hfill
%   \begin{subfigure}{0.48\linewidth}
%     \includegraphics[width=\linewidth]{figures/kmeans_synthetic_result.pdf}
%     \caption{$k$-means}
%   \end{subfigure}
%   \caption{Phân cụm trên dữ liệu tổng hợp 2D}
%   \label{fig:synthetic_result}
% \end{figure}

\subsubsection{Phân tích hội tụ}
% Gợi ý: vẽ đường cong log-likelihood qua các vòng lặp EM
% \begin{figure}[H]
%   \centering
%   \includegraphics[width=0.7\linewidth]{figures/convergence_curve.pdf}
%   \caption{Log-likelihood theo số vòng lặp EM}
%   \label{fig:convergence}
% \end{figure}
% TODO: Mô tả và phân tích đường cong hội tụ


% -----------------------------------------------------------
\subsection{Dữ liệu thực tế}
% -----------------------------------------------------------
% Gợi ý: Iris dataset, MNIST (2D qua PCA), Old Faithful, …
% TODO: Mô tả bộ dữ liệu

\subsubsection{Mô tả bộ dữ liệu}
% TODO: Viết nội dung ở đây

\subsubsection{Kết quả và thảo luận}
% TODO: Viết nội dung ở đây


% -----------------------------------------------------------
\subsection{Đánh giá định lượng}
% -----------------------------------------------------------
% Gợi ý nội dung:
%   - Adjusted Rand Index (ARI), Normalized Mutual Information (NMI)
%   - Silhouette Score (khi không có nhãn thật)
%   - BIC/AIC để chọn K

\begin{table}[H]
  \centering
  \caption{Kết quả đánh giá định lượng}
  \label{tab:quantitative}
  \begin{tabular}{lcccc}
    \toprule
    \textbf{Bộ dữ liệu} & \textbf{Thuật toán} & \textbf{ARI} & \textbf{NMI} & \textbf{Silhouette} \\
    \midrule
    % TODO: Điền kết quả vào đây
    Tổng hợp 2D & GMM-EM  & -- & -- & -- \\
    Tổng hợp 2D & $k$-means & -- & -- & -- \\
    \midrule
    Iris        & GMM-EM  & -- & -- & -- \\
    Iris        & $k$-means & -- & -- & -- \\
    \bottomrule
  \end{tabular}
\end{table}

% TODO: Nhận xét bảng kết quả
